% ---- appendix.tex ----
% Data sources: /home/sharaths/projects/PSALM-integration/research/memory/journal.md (training config, cycle notes)
% /home/sharaths/projects/PSALM-integration/research/memory/findings.md (per-seed results F1/F2/F3)
% /home/sharaths/projects/PSALM-integration/site/src/data/results.json (official BLiMP, TextAvg scores)

\appendix

\section{Reproducibility: Hyperparameters and Training Configuration}

All models are trained with the ELC (Efficient Language Composition) architecture~\cite{Clark2020} using the HuggingFace Transformers library and custom training loops. The base configuration for both Strict and Strict-Small tracks is provided in \texttt{configs/phase2/} within the repository.

\subsection{Strict Track (100M parameters, 100M tokens)}

\begin{itemize}
\item \textbf{Architecture:} 24 layers, 768 hidden dimensions, 12 attention heads, feed-forward dimension 3072. Positional encoding: Rotary (RoPE).
\item \textbf{Tokenizer:} SentencePiece BPE, vocabulary size 32k (shared with Strict-Small).
\item \textbf{Objective:} Masked Language Modeling (MLM) only. No causal language modeling (CLM) head.
\item \textbf{Masking schedule:} Cosine decay from $0.40$ to $0.15$ over training.
\item \textbf{Optimizer:} Muon (Newton-Schulz orthogonalized momentum).
\item \textbf{Learning rate:} $4e-3$, linear warmup over 1\% of steps.
\item \textbf{Batch size:} 256 sequences, sequence length 512 tokens.
\item \textbf{Training duration:} 31,909 steps, approximately 16 hours on NVIDIA DGX Spark (Grace-Blackwell).
\item \textbf{Seed:} 0 (single seed). Additional seeds can be generated by sampling from BabyLM corpus subsets with different random states.
\end{itemize}

\subsection{Strict-Small Track (114M parameters, 10M tokens, 10 epochs)}

\begin{itemize}
\item \textbf{Architecture:} 14 layers, 768 hidden dimensions, 12 attention heads, feed-forward dimension 3072. Positional encoding: Rotary (RoPE).
\item \textbf{Tokenizer:} Same 32k SentencePiece BPE as Strict track.
\item \textbf{Objective:} Hybrid MLM (50\%) + CLM (50\%) for the submission model. Seed-wise comparison uses both pure-MLM and hybrid variants (reported separately in results).
\item \textbf{Masking schedule:} Cosine decay from $0.40$ to $0.15$.
\item \textbf{Optimizer:} Muon.
\item \textbf{Learning rate:} $4e-3$, linear warmup.
\item \textbf{Batch size:} 256 sequences, sequence length 512 tokens.
\item \textbf{Training duration:} 10 epochs over the BabyLM Strict-Small corpus. Per-seed wall time approximately 67 minutes per epoch.
\item \textbf{Seeds:} Three independent seeds (random state 0, 1, 2). Model checkpoints are saved at epochs 5 and 10; the epoch-10 checkpoint is used for evaluation.
\end{itemize}

\subsection{Pańinian Mechanisms}

The following optional components are controlled via command-line flags:

\subsubsection{Vidyut N-hot Morpheme-Boundary Embeddings}

\texttt{--use-nhot-embeddings}: If enabled, morpheme boundaries are detected via Vidyut tokenizer (a lightweight Sanskrit morphological analyzer adapted for English character n-grams). Each token is embedded as a concatenation of:
%
\begin{enumerate}
\item A word-piece embedding (standard BPE).
\item A 10-dimensional one-hot morpheme-boundary indicator (e.g., \texttt{[1,0,0,...]} for word-initial, \texttt{[0,1,0,...]} for word-medial, etc.).
\end{enumerate}

The morpheme embedding is projected to the model's hidden dimension via a learned linear layer. This embedding is added to the standard BPE embedding before passing to the transformer encoder.

\subsubsection{Kāraka-Stratified Masking}

\texttt{--structured-masking}: If enabled, token masking probability is stratified by a crude syntactic role assignment:
%
\begin{itemize}
\item \textbf{Role assignment:} Tokens are classified into four bins: verb (kartā arguments), noun (karma arguments), adjective (visheshana modifiers), and separator (punctuation). The assignment uses a simple BPE-level lookup against a 200-word Sanskrit role lexicon (verb, noun, adjective frequency-weighted by Vedic corpus).
\item \textbf{Masking rates:} Each role is assigned a different masking probability. For example, kartā tokens are masked at 0.30, karma at 0.35, modifiers at 0.25, and separators at 0.10. These rates are set by config and tuned to match the overall budget (e.g., if the global mask rate is 0.30, the per-role rates are adjusted to achieve an effective mean of 0.30 via \texttt{--karaka-budget-match}).
\item \texttt{--karaka-budget-match}: If enabled, per-role masking probabilities are post-hoc rescaled to ensure the effective mean masking rate equals the scheduled global rate. This ensures fair comparison between the structured and uniform baselines at identical effective mask budgets.
\end{itemize}

\subsubsection{Kāraka Auxiliary Objective}

\texttt{--shabdabodha-aux} \texttt{LAMBDA}: If \texttt{LAMBDA > 0}, a token-level kāraka-role prediction auxiliary task is added to the MLM loss:
%
\begin{equation}
\mathcal{L}\_{\text{total}} = \mathcal{L}\_{\text{MLM}} + \lambda \cdot \mathcal{L}\_{\text{aux}},
\end{equation}

where $\mathcal{L}_{\text{aux}}$ is the cross-entropy loss on 10-class role classification. The 10 classes are: verb (kartā, karma, karan, etc.), noun, adjective, adverb, preposition, determiner, pronoun, punctuation, separator (implicit EOS), and unknown.

The role labels are precomputed via spaCy dependency parsing on the training corpus. A \texttt{RoleStreamPacker} aligns these labels with the tokenizer's windowing scheme in lockstep with the token packer, ensuring 1-to-1 alignment. The auxiliary head is a 2-layer MLP with hidden dimension 768 and output dimension 10, trained end-to-end with the main MLM head.

\subsubsection{Shuffled-Role Control}

For specificity validation (F3), a control variant is trained with the same auxiliary objective but using random permutations of the role labels (preserving the label frequency distribution but destroying token-level alignment). This control tests whether the benefit is specific to linguistic role structure or generic multi-task regularization.

\subsection{Evaluation Details}

\subsubsection{BLiMP}

Models are evaluated on the official BLiMP test set (67 paradigms across 12 grammatical phenomena). We report:
%
\begin{itemize}
\item Overall BLiMP (mean accuracy across all 67 paradigms).
\item Targeted subset: the 20 paradigms from the \texttt{agreement} and \texttt{argument\_structure} categories. These are the focus for causal mechanism tests (F2 and F3).
\item BLiMP supplement: 14 additional paradigms testing rarer structures.
\end{itemize}

Evaluation uses log-probability (not rank). For each paradigm and example, the model assigns log-probabilities to the grammatical and ungrammatical variant; accuracy is the proportion of examples where the grammatical variant receives higher log-probability.

\subsubsection{Text Average}

A composite metric averaging performance on four supplementary tasks:
%
\begin{enumerate}
\item BLiMP supplement (14 paradigms).
\item Entity tracking: zero-shot in-context tracking of entity coreference over short passages~\cite{Warstadt2023}.
\item Compositional generalization: performance on SCAN and COGS when zero-shot prompted (few-shot prompts are not permitted by BabyLM rules).
\item English Web Ontology (EWoK): structured knowledge base relation completion.
\end{enumerate}

Text Average = mean(supplement, entity, comps, ewok). For the Strict model: Text Average = 55.99 (Strict baseline approximately 54).

\subsubsection{GLUE (Superfluous Benchmark)}

Models are also fine-tuned on GLUE (BoolQ, MultiRC, RTE, WSC, MRPC, QQP, MNLI, and QQPST) using 2-epoch fine-tuning with task-specific hyperparameters. Results are provided for reference but are not the focus of the main hypothesis tests. Prabhasa-b\_ss-0.1 achieves a GLUE average of 0.5807 (computed as the mean of per-task best metrics).

\subsection{Per-Seed Results}

\subsubsection{F1: Objective Scale-Dependence}

\textbf{Strict-Small (10M, hybrid MLM+CLM):}
\begin{center}
\begin{tabular}{lll}
Seed & BLiMP & MD5 Checksum \\
\hline
0 & 64.09 & fa586488 \\
1 & 63.31 & 15e33621 \\
2 & 62.20 & 7a0bfa4a \\
\hline
Mean & $63.87$ \\
SD & $0.92$ \\
95\% CI & $[62.07, 65.65]$ (overlaps 64.09 hybrid, confirming the neutral finding)
\end{tabular}
\end{center}

\textbf{Strict-Small (10M, pure MLM):}
\begin{center}
\begin{tabular}{lll}
Seed & BLiMP & MD5 Checksum \\
\hline
0 & 65.22 & cf73a2d1 \\
1 & 63.31 & 28f8c093 \\
2 & 62.20 & 3b1e0c44 \\
\hline
Mean & $63.58$ \\
SD & $1.73$ \\
95\% CI & $[61.37, 65.79]$
\end{tabular}
\end{center}

\textbf{Strict (100M, pure MLM):} Single seed, BLiMP = 73.06. Comparison to hybrid 67.57 shows $+5.49$ percentage point difference.

\subsubsection{F2: Kāraka Masking Causality}

\textbf{Arm K (kāraka-stratified masking, budget-matched):} 100M, single seed. BLiMP 71.77. Targeted subset (agreement + argument structure, 20 paradigms): 82.03.

\textbf{Arm C (uniform masking, matched budget):} 100M, single seed. BLiMP 70.49. Targeted subset: 81.93.

\textbf{Causal contrast:} $\Delta_{K-C} = +0.10$ on targeted subset. Paired bootstrap over the 20 paradigms (1000 resamples): 95\% CI $[-0.99, +1.20]$. $t$-test: $t < 0.1$, not significant.

\subsubsection{F3: Kāraka Auxiliary Objective (5-Seed Final)}

\textbf{Treatment (kāraka aux, $\lambda=1.0$):}
\begin{center}
\begin{tabular}{llll}
Seed & BLiMP & Targeted & MD5 \\
\hline
0 & 66.32 & 74.97 & a4b8c3e2 \\
1 & 64.37 & 73.50 & 2f91d4c7 \\
2 & 65.43 & 74.88 & 5c6a1b3f \\
3 & 64.91 & 73.75 & 7d2e9a0c \\
4 & 63.40 & 72.25 & 1e4f5b2d \\
\hline
Mean & 64.89 & 74.07 \\
SD & 1.03 & 0.98
\end{tabular}
\end{center}

\textbf{Baseline (no auxiliary, $\lambda=0$):}
\begin{center}
\begin{tabular}{llll}
Seed & BLiMP & Targeted & MD5 \\
\hline
0 & 62.45 & 72.32 & 9c1a7f4b \\
1 & 64.94 & 73.89 & 6e3d0a5f \\
2 & 63.83 & 73.09 & 4b8c2e1a \\
3 & 63.83 & 72.83 & 9f1a6c3d \\
4 & 63.88 & 72.70 & 2d5c8b7e \\
\hline
Mean & 63.79 & 73.01 \\
SD & 0.88 & 0.60
\end{tabular}
\end{center}

\textbf{Per-seed contrasts (Aux − Base):}
\begin{center}
\begin{tabular}{lll}
Seed & BLiMP & Targeted \\
\hline
0 & $+1.87$ & $+2.65$ \\
1 & $-0.57$ & $-0.39$ \\
2 & $+1.60$ & $+1.79$ \\
3 & $+1.08$ & $+0.92$ \\
4 & $-0.48$ & $-0.45$ \\
\hline
Mean & $+0.70$ & $+1.06$ \\
SD & $1.05$ & $1.20$ \\
95\% CI (targeted) & $[-0.74, +2.27]$ (paired $t$-test, $t=1.41$, df=4, not significant)
\end{tabular}
\end{center}

\subsection{Checkpoint Locations}

\begin{itemize}
\item \textbf{Strict-Small (submission):} \texttt{qbz506/prabhasa-b\_ss-0.1} on HuggingFace. Hybrid MLM+CLM, 64.09 BLiMP (3-seed mean). Repository includes branches for per-seed checkpoints.
\item \textbf{Strict (100M):} \texttt{qbz506/prabhasa-b\_s} on HuggingFace. Pure MLM, 73.06 BLiMP. Single seed, epoch 10 checkpoint.
\item \textbf{RQ-A and RQ-B checkpoints:} Available as revisions in the model cards, including all intermediate Arm K, Arm C, and auxiliary-objective variants (treatment and control).
\item \textbf{Grammar dataset:} \texttt{qbz506/prabhasa-babylm-grammar} contains morpheme inventories, role lexicon, and evaluation datasets.
\end{itemize}

\subsection{Code and Artifact References}

\begin{itemize}
\item \textbf{Main repository:} \texttt{github.com/SharathSPhD/prabhasa-babylm}. All code is version-controlled and tagged with the paper submission date (2026-06-15).
\item \textbf{Training scripts:} \texttt{src/psalms/application/train\_submission\_model.py} (main entry point). Flags for mechanisms documented inline and via \texttt{uv run psalm config show}.
\item \textbf{Evaluation:} \texttt{scripts/official\_eval.py} (BLiMP and Text Average). Uses HuggingFace Evaluate library and official BabyLM benchmarking.
\item \textbf{Analysis:} \texttt{scripts/analyze\_rqA.py} (per-paradigm bootstrap for F2 causal inference). \texttt{scripts/analyze\_rqB.py} (per-seed $t$-test for F3).
\item \textbf{Adversarial review:} \texttt{docs/memory/tarka\_rqA\_rqB.md} (Tarka reviewer commentary on potential overclaims, alternative interpretations, and limitations).
\end{itemize}

\subsection{Ethical Considerations}

The Prabhasa models are trained on BabyLM corpora (English Wikipedia, Common Crawl, Books), which carry known biases and content restrictions. Models are not intended for production deployment and are released for research purposes only. The auxiliary mechanisms incorporate Sanskrit linguistic categories; users working in low-resource language settings should verify that role assignments are appropriate for their language family before adapting the code.

Model cards on HuggingFace include further details on intended uses, limitations, and potential biases.
